\input{common/onetextpage.tex}
\newcounter{nappendix}
\renewcommand{\thenappendix}{\Alph{nappendix}}
\newcommand{\newappendix}[1]{%
    \refstepcounter{nappendix}
    \onetextpage{Appendix \thenappendix}
    \chapter*{Appendix \thenappendix: #1}
    \label{appendix:\thenappendix}
    \addcontentsline{toc}{chapter}{Appendix \thenappendix: #1}
}

%% Edit past this comment line.

\newappendix{RIASEC-to-Topic Cluster Mapping}

\setcounter{table}{0}
\renewcommand{\thetable}{A.\arabic{table}}

Table~\ref{tab:riasec-mapping} documents the mapping between Holland's six RIASEC types and the eight topic clusters presented in the interest discovery interface. Each topic tile is assigned a primary RIASEC type and a secondary RIASEC type with associated weights. When a student selects a tile, both weights are added to the corresponding dimensions of the 6-dimensional RIASEC vector. After all five elicitation steps (tile selection, pairwise comparisons, scenario confirmation, dealbreaker filter, and confidence weighting), the vector is L2-normalised before use in retrieval and fit score computation.

\begin{table}[htbp]
\centering
\small
\renewcommand{\arraystretch}{1.4}
\caption{RIASEC-to-Topic Cluster Weight Mapping}
\label{tab:riasec-mapping}
\begin{tabularx}{\textwidth}{|p{3.5cm}|p{2.5cm}|p{1.2cm}|p{3cm}|p{1.2cm}|}
\hline
\textbf{Topic Cluster} & \textbf{Primary RIASEC} & \textbf{Weight} & \textbf{Secondary RIASEC} & \textbf{Weight} \\
\hline
Computers \& Digital & Investigative (I) & 0.8 & Realistic (R) & 0.5 \\
\hline
Science \& Research & Investigative (I) & 0.9 & Realistic (R) & 0.4 \\
\hline
Engineering \& Building & Realistic (R) & 0.9 & Investigative (I) & 0.5 \\
\hline
Arts \& Design & Artistic (A) & 0.9 & Enterprising (E) & 0.4 \\
\hline
Business \& Leadership & Enterprising (E) & 0.9 & Social (S) & 0.4 \\
\hline
People \& Society & Social (S) & 0.9 & Enterprising (E) & 0.4 \\
\hline
Nature \& Environment & Realistic (R) & 0.8 & Investigative (I) & 0.4 \\
\hline
Organization \& Data & Conventional (C) & 0.9 & Investigative (I) & 0.4 \\
\hline
\end{tabularx}
\end{table}

\noindent\textit{Note: Weights reflect the relative centrality of each RIASEC dimension to the topic cluster as assessed by the development team and reviewed by the faculty advisor. These values are intended as a reasonable initialisation and may be refined through user testing.}

%% Remove the following appendix before publishing your work.

\newappendix{About \LaTeX}
LaTeX (stylized as \LaTeX) is a software system for typesetting documents. LaTeX markup describes the content and layout of the document,
as opposed to the formatted text found in WYSIWYG word processors like Google Docs, LibreOffice Writer, and Microsoft Word.
The writer uses markup tagging conventions to define the general structure of a document,
to stylize text throughout a document (such as bold and italics), and to add citations and cross-references.

LaTeX is widely used in academia for the communication and publication of scientific documents and technical note-taking in many fields,
owing partially to its support for complex mathematical notation. It also has a prominent role in the preparation and publication of books and articles
that contain complex multilingual materials, such as Arabic and Greek.

Overleaf has also provided a 30-minute guide on how you can get started on using \LaTeX. \cite{overleaf:learnlatex}